\documentclass[11pt]{article}
\usepackage[margin=1in]{geometry}
\usepackage{amsmath,amssymb}
\usepackage{booktabs}
\usepackage[colorlinks=true,linkcolor=blue,citecolor=blue]{hyperref}

\newcommand{\R}{\mathbb{R}}
\newcommand{\E}{\mathbb{E}}
\newcommand{\bs}[1]{\boldsymbol{#1}}
\newcommand{\norm}[1]{\left\lVert #1 \right\rVert}

\title{An outcome surrogate for granular interactions:\\
proper orthogonal decomposition with Gaussian process regression}
\author{}
\date{26 August 2026}

\begin{document}
\maketitle

\begin{abstract}
We describe the construction of a parametric surrogate for the outcome of a granular
interaction computed by a three-dimensional particle dynamics simulation. The particle state
is interpolated onto a coarse Eulerian grid, the resulting field ensemble is reduced by
proper orthogonal decomposition (POD), and the map from interaction parameters to modal
coefficients is regressed by Gaussian process (GP) regression. The purpose of the document
is to fix notation and dimensions; measured results from a pilot ensemble are reported in
Section~\ref{sec:measured}.
\end{abstract}

\section{Problem statement and notation}
\label{sec:notation}

Let $d = 3$ denote the spatial dimension and $\Omega \subset \R^{d}$ the simulation domain,
a rectangular box. The full-order model (FOM) advances $n_p$ particles subject to gravity,
contact forces and boundary contacts. Particle $i \in \{1, \ldots, n_p\}$ carries a position
$\bs{x}_i \in \R^{d}$, a velocity $\bs{v}_i \in \R^{d}$ and a mass $m_i > 0$. Nothing below
depends on the particular force law, and the formulation applies unchanged if the FOM is
later extended with rigid bodies or continuum fields.

The state is stored with the degrees of freedom of each particle contiguous, position
followed by velocity:
\begin{equation}
  \bs{q}(t) \in \R^{n_q}, \quad n_q = 2 d n_p,
  \qquad
  \bs{q} = \big[\, \bs{x}_1, \bs{v}_1, \; \bs{x}_2, \bs{v}_2, \; \ldots, \;
  \bs{x}_{n_p}, \bs{v}_{n_p} \,\big],
  \label{eq:state}
\end{equation}
so that, with components indexed from $1$,
\begin{equation}
  q_{2d(i-1) + j} = (\bs{x}_i)_j, \qquad
  q_{2d(i-1) + d + j} = (\bs{v}_i)_j, \qquad j = 1, \ldots, d .
  \label{eq:state-index}
\end{equation}
For $d = 3$ the stride between consecutive particles is $2d = 6$. This is the layout used by
the implementation; it is recorded here because every index expression in
Section~\ref{sec:encoder} inherits it. For the ensemble reported below $n_p = 3375$, hence
$n_q = 20\,250$.

Each realization of the FOM is determined by a parameter vector
\begin{equation}
  \bs{\theta} \in \mathcal{P} \subset \R^{p}, \qquad p = 7,
  \label{eq:theta}
\end{equation}
where $\mathcal{P}$ is a product of intervals. The components of $\bs{\theta}$ are the
$d = 3$ components of the initial velocity applied to every particle, the release height,
the per-particle mass, the coefficient of restitution $e \in [0,1]$, and the Coulomb
friction coefficient $\mu \geq 0$. The first four are \emph{actions}, in the sense that an
agent acting in the environment would select them; the last three are \emph{material}
parameters, which system identification would instead infer from observation.

We seek a surrogate for a coarse-grained functional of the terminal state, not for
$\bs{q}$ itself. This is deliberate: $\bs{q}$ is chaotic, whereas coarse-grained
functionals of a granular assembly are reproducible.

\section{Interpolation onto the grid}
\label{sec:encoder}

Let the grid have $N_j$ nodes along axis $j$, covering $\Omega$ inclusively, so the spacing
is $h_j = (\Omega_j^{\max} - \Omega_j^{\min})/(N_j - 1)$. Nodes are labelled by a single
index
\begin{equation}
  \alpha \in \{1, \ldots, N\}, \qquad N = N_1 N_2 N_3 ,
  \label{eq:node-index}
\end{equation}
obtained from the triple $(\alpha_1, \alpha_2, \alpha_3)$ of per-axis indices by
$\alpha = (\alpha_1 - 1) N_2 N_3 + (\alpha_2 - 1) N_3 + \alpha_3$. Presently
$(N_1, N_2, N_3) = (32, 32, 16)$ and $N = 16\,384$. Let $\bs{y}_\alpha \in \R^{d}$ denote
the position of node $\alpha$.

Define $n_c = 1 + d = 4$ \emph{channels}: the mass density and the $d$ components of the
momentum density. Particle quantities are interpolated onto the grid with the product of
one-dimensional hat functions,
\begin{equation}
  w_{i\alpha} = \prod_{j=1}^{d} \max\!\left( 0, \; 1 -
  \frac{\left| (\bs{x}_i)_j - (\bs{y}_\alpha)_j \right|}{h_j} \right) \; \geq 0 ,
  \label{eq:weights}
\end{equation}
which are nonzero on the $2^d = 8$ nodes surrounding $\bs{x}_i$ and form a partition of
unity,
\begin{equation}
  \sum_{\alpha=1}^{N} w_{i\alpha} = 1 \qquad \text{for every } i .
  \label{eq:partition}
\end{equation}
The interpolation is then
\begin{equation}
  \rho_\alpha = \sum_{i=1}^{n_p} m_i w_{i\alpha},
  \qquad
  (\rho u_j)_\alpha = \sum_{i=1}^{n_p} m_i (\bs{v}_i)_j \, w_{i\alpha},
  \qquad j = 1, \ldots, d,
  \label{eq:deposit}
\end{equation}
and \eqref{eq:partition} gives exact conservation of mass,
\begin{equation}
  \sum_{\alpha=1}^{N} \rho_\alpha
  = \sum_{i=1}^{n_p} m_i \sum_{\alpha=1}^{N} w_{i\alpha}
  = \sum_{i=1}^{n_p} m_i .
  \label{eq:conservation}
\end{equation}
Assigning each particle wholly to its nearest node also satisfies \eqref{eq:conservation}
but is discontinuous in $\bs{x}_i$; the weights \eqref{eq:weights} are continuous, which is
required for the ensemble of Section~\ref{sec:pod} to depend continuously on
$\bs{\theta}$.

Write $\bs{s}_c \in \R^{N}$ for the field of channel $c \in \{1, \ldots, n_c\}$, with
$c = 1$ the mass channel and $c = 1 + j$ the $j$th momentum component, and
$\bs{s} \in \R^{n_c N}$ for their concatenation, $\dim \bs{s} = 65\,536$. The map
$\bs{q} \mapsto \bs{s}$ defined by \eqref{eq:deposit} is denoted
$\mathcal{E} : \R^{n_q} \to \R^{n_c N}$.

The momentum channels are retained because $\bs{s}_1$ alone is not a Markov state: two
assemblies with identical mass fields and different velocity fields evolve differently.
They are, however, informative only while the material is in motion; see
Section~\ref{sec:measured}.

\section{Reduction: proper orthogonal decomposition}
\label{sec:pod}

\subsection{Snapshot ensemble}

Let $\{\bs{\theta}^{(m)}\}_{m=1}^{M}$ be a design of experiments in $\mathcal{P}$, obtained
here by a scrambled Sobol sequence, a quasi-Monte Carlo (QMC) construction whose marginal
projections are uniform by construction. For each $m$ the FOM is advanced to a terminal time
$T$ and interpolated by \eqref{eq:deposit}, giving snapshots $\bs{s}_c^{(m)} \in \R^{N}$.

For each channel $c$ independently, collect the $M$ snapshots as the rows of
\begin{equation}
  X_c \in \R^{M \times N},
  \qquad
  \left( X_c \right)_{m\alpha} = \left( \bs{s}_c^{(m)} \right)_{\!\alpha} ,
  \label{eq:snapshot-matrix}
\end{equation}
that is, entry $(m, \alpha)$ of $X_c$ is the value of channel $c$ at grid node $\alpha$
in realization $m$: the row index runs over realizations, the column index over grid nodes.
Subtract the ensemble mean field from every row,
\begin{equation}
  \tilde{X}_c = X_c - \bs{1}_M \bar{\bs{s}}_c^{\,\top},
  \qquad
  \bar{\bs{s}}_c = \frac{1}{M} \sum_{m=1}^{M} \bs{s}_c^{(m)} \in \R^{N},
  \label{eq:centering}
\end{equation}
with $\bs{1}_M \in \R^{M}$ the vector of ones, so that $\bs{1}_M \bar{\bs{s}}_c^{\,\top}
\in \R^{M \times N}$ repeats $\bar{\bs{s}}_c^{\,\top}$ in each of its $M$ rows.

A separate basis is constructed for each channel. A single basis spanning all $n_c$ channels
jointly would require a relative scaling between mass density and momentum density. Such a
scaling is not physically determined: $\rho \geq 0$ whereas $\rho u_j$ is sign-indefinite,
and the two have different units and substantially different dynamic ranges. Under any fixed
scaling the channel assigned the larger magnitude dominates the singular spectrum, and the
fraction of captured variance ceases to be a meaningful diagnostic. Independent bases
additionally permit channel-dependent truncation ranks.

\subsection{Singular value decomposition}

The singular value decomposition (SVD) of the centered snapshot matrix is
\begin{equation}
  \tilde{X}_c = U_c \Sigma_c V_c^{\top},
  \qquad
  U_c \in \R^{M \times r_c}, \quad
  \Sigma_c \in \R^{r_c \times r_c}, \quad
  V_c \in \R^{N \times r_c},
  \label{eq:svd}
\end{equation}
with $U_c^{\top} U_c = V_c^{\top} V_c = I_{r_c}$, $\Sigma_c = \mathrm{diag}(\sigma_{c,1},
\ldots, \sigma_{c,r_c})$, $\sigma_{c,1} \geq \cdots \geq \sigma_{c,r_c} > 0$, and
$r_c = \mathrm{rank}\,\tilde{X}_c \leq \min(M-1, N)$. The bound $M-1$ rather than $M$ follows
from \eqref{eq:centering}, which removes one degree of freedom from the row space. For the
present ensemble $M = 150 \ll N$, so $r_c \leq 149$ and the SVD is computed in economy form
at cost $O(M^2 N)$.

The columns $\bs{\phi}_{c,k} \in \R^{N}$ of $V_c$ are the POD modes of channel $c$:
orthonormal spatial fields on the grid. Retaining the leading $K_c \leq r_c$ modes defines
\begin{equation}
  \Phi_c = \left[ \bs{\phi}_{c,1}, \ldots, \bs{\phi}_{c,K_c} \right] \in \R^{N \times K_c},
  \qquad \Phi_c^{\top} \Phi_c = I_{K_c},
  \label{eq:basis}
\end{equation}
and the modal coefficients of a snapshot follow by orthogonal projection,
\begin{equation}
  \bs{a}_c^{(m)} = \Phi_c^{\top} \left( \bs{s}_c^{(m)} - \bar{\bs{s}}_c \right)
  \in \R^{K_c},
  \qquad
  \bs{s}_c^{(m)} \approx \bar{\bs{s}}_c + \Phi_c \bs{a}_c^{(m)} .
  \label{eq:projection}
\end{equation}
The truncation rank $K_c$ is chosen per channel; it is not assumed equal across $c$.

Two properties motivate this basis. First, by the Eckart--Young theorem the truncation
\eqref{eq:basis} is optimal: among all rank-$K_c$ matrices,
\begin{equation}
  \min_{\mathrm{rank}\,Y \leq K_c} \norm{\tilde{X}_c - Y}_F^2
  = \sum_{k = K_c+1}^{r_c} \sigma_{c,k}^2 ,
  \label{eq:eckart-young}
\end{equation}
where $\norm{\cdot}_F$ is the Frobenius norm. No other $K_c$-dimensional linear subspace
represents the ensemble with smaller squared error. Second, \eqref{eq:eckart-young} makes the
truncation error available from the spectrum alone, before any regression is attempted.

\subsection{Retained variance, and the term ``energy''}
\label{sec:energy}

Because
\begin{equation}
  \norm{\tilde{X}_c}_F^2
  = \sum_{k=1}^{r_c} \sigma_{c,k}^2
  = \sum_{m=1}^{M} \norm{ \bs{s}_c^{(m)} - \bar{\bs{s}}_c }_2^2 ,
  \label{eq:frobenius-identity}
\end{equation}
the natural truncation diagnostic is the fraction of that total carried by the retained
modes,
\begin{equation}
  \eta_c(K) = \frac{\sum_{k=1}^{K} \sigma_{c,k}^2}{\sum_{k=1}^{r_c} \sigma_{c,k}^2}
  \; = \; 1 - \frac{\min_{\mathrm{rank}\,Y \leq K} \norm{\tilde{X}_c - Y}_F^2}
  {\norm{\tilde{X}_c}_F^2} \; \in [0, 1] .
  \label{eq:energy}
\end{equation}
By \eqref{eq:frobenius-identity} the denominator is $(M-1)$ times the sample variance of the
field summed over all $N$ nodes, so $\eta_c$ is exactly the \emph{fraction of ensemble
variance retained}.

The quantity commonly reported as ``retained energy'', and quoted as a percentage, is the
same $\eta_c(K)$ of \eqref{eq:energy}; the two names refer to one quantity and there is no
separate formula. The name is literal only in the setting where POD was first popularized,
in which the snapshots are velocity fields: then $\sum_k \sigma_k^2$ is proportional to the
total fluctuating kinetic energy of the ensemble, and truncation discards a definite amount
of it. For a mass-density channel no energy is involved and the term is an analogy resting
only on $\eta_c$ being a normalized sum of squares. We therefore report $\eta_c$ as retained
variance, and note the synonym once, here.

$\eta_c$ is a statement about the ensemble used to build the basis. It is not a
generalization estimate, and Section~\ref{sec:errors} does not use it as one.

\subsection{Error measures}
\label{sec:errors}

All errors are measured between fields on the grid, in $\R^{N}$: the reference is the
interpolated field $\bs{s}_c^{(m)}$ of \eqref{eq:deposit}, not the particle state
$\bs{q}$. This is by construction rather than by convenience. The grid field is the declared
quantity of interest, so the information discarded by $\mathcal{E}$ is part of the problem
statement and not an error to be charged against the reduction or the regression. A metric
comparing predictions against $\bs{q}$ would instead measure the coarse-graining, which is a
modeling choice made in Section~\ref{sec:notation}.

Let $\mathcal{T}$ be a set of realizations held out from the construction of
$\bar{\bs{s}}_c$, of $\Phi_c$, and of the regressor of Section~\ref{sec:gp}. Two errors are
reported separately. The first replaces the true coefficients into
\eqref{eq:projection} and so isolates the truncation:
\begin{equation}
  \varepsilon^{\text{proj}}_c(K_c) = \left(
  \frac{\sum_{m \in \mathcal{T}} \norm{
    \big[ \bar{\bs{s}}_c + \Phi_c \Phi_c^{\top}
    ( \bs{s}_c^{(m)} - \bar{\bs{s}}_c ) \big] - \bs{s}_c^{(m)} }_2^2 }{
  \sum_{m \in \mathcal{T}} \norm{ \bs{s}_c^{(m)} }_2^2 } \right)^{1/2} ,
  \label{eq:error-proj}
\end{equation}
in which the bracketed term is the orthogonal projection of the true field onto the affine
subspace $\bar{\bs{s}}_c + \mathrm{range}\,\Phi_c$. The second uses the regressed
coefficients $\bs{m}_c(\bs{\theta}^{(m)})$ of \eqref{eq:gp-mean} and so measures the
deployed surrogate:
\begin{equation}
  \varepsilon^{\text{sur}}_c(K_c) = \left(
  \frac{\sum_{m \in \mathcal{T}} \norm{
    \big[ \bar{\bs{s}}_c + \Phi_c \bs{m}_c(\bs{\theta}^{(m)}) \big]
    - \bs{s}_c^{(m)} }_2^2 }{
  \sum_{m \in \mathcal{T}} \norm{ \bs{s}_c^{(m)} }_2^2 } \right)^{1/2}
  \; \geq \; \varepsilon^{\text{proj}}_c(K_c) ,
  \label{eq:error-sur}
\end{equation}
the inequality holding because the projection is the minimizer over coefficients. The gap
between \eqref{eq:error-proj} and \eqref{eq:error-sur} is attributable to the regression
alone, so the two remedies -- enlarging $K_c$ and improving the regressor -- can be
prioritized from measurement rather than assumption.

Three choices in \eqref{eq:error-proj}--\eqref{eq:error-sur} are deliberate. The norms are
taken on the \emph{uncentered} field, because that is what the surrogate must produce;
reporting the error of $\bs{s}_c - \bar{\bs{s}}_c$ instead would credit the method with the
mean field, which is free. The ratio is formed once over the whole of $\mathcal{T}$ rather
than by averaging per-realization ratios, so that realizations of small total mass, whose
individual ratios are ill-conditioned, do not dominate. And $\mathcal{T}$ is excluded from
$\bar{\bs{s}}_c$ and $\Phi_c$: computing either from all $M$ realizations leaks the test set
into the basis and understates \eqref{eq:error-proj}.

\section{Regression: Gaussian process on modal coefficients}
\label{sec:gp}

Composing \eqref{eq:projection} with the parameter-to-state map reduces the task to the
regression of $K_c$ scalar functions of $\bs{\theta}$,
\begin{equation}
  f_{c,k} : \mathcal{P} \subset \R^{p} \to \R,
  \qquad
  f_{c,k}\big(\bs{\theta}^{(m)}\big) = \big( \bs{a}_c^{(m)} \big)_k,
  \qquad k = 1, \ldots, K_c .
  \label{eq:regression-task}
\end{equation}
Each $f_{c,k}$ is modeled as an independent GP: a prior over functions under which any
finite collection of function values is jointly Gaussian, specified by a mean function, taken
zero after centering, and a covariance (kernel) function $\kappa : \mathcal{P} \times
\mathcal{P} \to \R$.

We adopt the Mat\'ern kernel with smoothness parameter $\nu = 5/2$ and automatic relevance
determination (ARD), that is, one length scale per input,
\begin{equation}
  \kappa(\bs{\theta}, \bs{\theta}') = \sigma_f^2
  \left( 1 + \sqrt{5}\, \tau + \tfrac{5}{3} \tau^2 \right) e^{-\sqrt{5}\,\tau},
  \qquad
  \tau = \left( \sum_{i=1}^{p} \frac{(\theta_i - \theta_i')^2}{\ell_i^2} \right)^{1/2},
  \label{eq:matern}
\end{equation}
with signal variance $\sigma_f^2 > 0$ and length scales $\ell_i > 0$. Sample paths of
\eqref{eq:matern} are twice differentiable, a weaker and, for a system with contact, more
defensible assumption than the analyticity implied by a squared-exponential kernel.
Observations are modeled as $y^{(m)} = f_{c,k}(\bs{\theta}^{(m)}) + \epsilon^{(m)}$ with
$\epsilon^{(m)} \sim \mathcal{N}(0, \sigma_n^2)$ independent. Here $\sigma_n^2$ absorbs the
irreducible scatter of a chaotic system and is therefore a physically meaningful quantity
rather than a regularization device.

Let $\Theta \in \R^{M_{\mathrm{tr}} \times p}$ collect the $M_{\mathrm{tr}}$ training
parameter vectors, $\bs{y} \in \R^{M_{\mathrm{tr}}}$ the corresponding coefficients for a
fixed pair $(c,k)$, $\mathrm{K} \in \R^{M_{\mathrm{tr}} \times M_{\mathrm{tr}}}$ the Gram
matrix $\mathrm{K}_{mn} = \kappa(\bs{\theta}^{(m)}, \bs{\theta}^{(n)})$, and
$\bs{\kappa}_\ast \in \R^{M_{\mathrm{tr}}}$ the vector
$(\bs{\kappa}_\ast)_m = \kappa(\bs{\theta}^{(m)}, \bs{\theta}_\ast)$ at a query
$\bs{\theta}_\ast$. The posterior is Gaussian with mean and variance
\begin{align}
  m_{c,k}(\bs{\theta}_\ast) &= \bs{\kappa}_\ast^{\top}
    \left( \mathrm{K} + \sigma_n^2 I_{M_{\mathrm{tr}}} \right)^{-1} \bs{y},
  \label{eq:gp-mean} \\
  v_{c,k}(\bs{\theta}_\ast) &= \kappa(\bs{\theta}_\ast, \bs{\theta}_\ast)
    - \bs{\kappa}_\ast^{\top}
    \left( \mathrm{K} + \sigma_n^2 I_{M_{\mathrm{tr}}} \right)^{-1} \bs{\kappa}_\ast .
  \label{eq:gp-var}
\end{align}
Hyperparameters $(\sigma_f, \ell_1, \ldots, \ell_p, \sigma_n)$ are obtained by maximizing the
log marginal likelihood
\begin{equation}
  \log p(\bs{y} \mid \Theta) = -\tfrac{1}{2} \bs{y}^{\top} A^{-1} \bs{y}
  - \tfrac{1}{2} \log \det A - \tfrac{M_{\mathrm{tr}}}{2} \log 2\pi,
  \qquad A = \mathrm{K} + \sigma_n^2 I ,
  \label{eq:mll}
\end{equation}
at cost $O(M_{\mathrm{tr}}^3)$ per evaluation through the Cholesky factorization of $A$. For
$M_{\mathrm{tr}} = O(10^2$--$10^3)$ this is negligible.

Three properties are relevant here. First, \eqref{eq:gp-var} supplies a predictive variance
in closed form, which permits the model's stated uncertainty to be compared against the
measured scatter of the FOM. Second, exact inference is available at the sample sizes
affordable, where a high-capacity parametric regressor would be data-starved. Third, the
fitted length scales rank the sensitivity of the outcome to each component of $\bs{\theta}$,
since $\ell_i \to \infty$ removes input $i$ from \eqref{eq:matern}; this follows from
\eqref{eq:mll} without additional sampling.

\section{The composed surrogate and its dimensions}
\label{sec:dimensions}

The surrogate for channel $c$ is
\begin{equation}
  \hat{\bs{s}}_c(\bs{\theta}) = \bar{\bs{s}}_c + \Phi_c \bs{m}_c(\bs{\theta}),
  \qquad
  \bs{m}_c(\bs{\theta}) = \big[ m_{c,1}(\bs{\theta}), \ldots, m_{c,K_c}(\bs{\theta})
  \big]^{\top} \in \R^{K_c} .
  \label{eq:surrogate}
\end{equation}
Because \eqref{eq:surrogate} is affine in the coefficients, $\hat{\bs{s}}_c$ is the exact
conditional mean of the reduced representation,
$\E[\bs{s}_c \mid \bs{\theta}] = \bar{\bs{s}}_c + \Phi_c \E[\bs{a}_c \mid \bs{\theta}]$;
the independence assumed across $k$ therefore does \emph{not} bias the predicted mean field.
It affects only the predictive covariance, which under independence is the rank-$K_c$ matrix
\begin{equation}
  \mathrm{Cov}\!\left[ \bs{s}_c \mid \bs{\theta} \right]
  \approx \Phi_c \, \mathrm{diag}\!\left( v_{c,1}(\bs{\theta}), \ldots,
  v_{c,K_c}(\bs{\theta}) \right) \Phi_c^{\top} \in \R^{N \times N},
  \label{eq:field-cov}
\end{equation}
and which omits any variance in the discarded subspace.

The number of scalar GPs required is
\begin{equation}
  n_{\mathrm{GP}} = \sum_{c \in \mathcal{C}} K_c ,
  \label{eq:ngp}
\end{equation}
where $\mathcal{C} \subseteq \{1, \ldots, n_c\}$ is the set of predicted channels.
Predicting all channels at a common rank $K$ gives $n_{\mathrm{GP}} = n_c K = 4K$, and the
surrogate is a map $\R^{p} \to \R^{4K}$ followed by $n_c$ independent affine liftings
$\R^{K} \to \R^{N}$. For the terminal-state task of Section~\ref{sec:measured} we take
$\mathcal{C} = \{1\}$ and $K_1 = 20$, so $n_{\mathrm{GP}} = 20$ and the surrogate is
$\R^{7} \to \R^{20} \to \R^{16\,384}$. Table~\ref{tab:dims} collects the dimensions.

\begin{table}[t]
\centering
\begin{tabular}{llr}
\toprule
Symbol & Meaning & Value \\
\midrule
$d$ & spatial dimension & 3 \\
$n_p$ & particles per realization & 3375 \\
$n_q$ & FOM state dimension, $2 d n_p$ & 20\,250 \\
$p$ & parameters, $\dim \bs{\theta}$ & 7 \\
$n_c$ & channels, $1 + d$ & 4 \\
$(N_1,N_2,N_3)$ & grid nodes per axis & $(32,32,16)$ \\
$N$ & nodes per channel, $N_1 N_2 N_3$ & 16\,384 \\
$\dim \bs{s}$ & interpolated state, $n_c N$ & 65\,536 \\
$M$ & realizations (snapshots) & 150 \\
$N_t$ & recorded instants per realization & 10 \\
$r_c$ & $\mathrm{rank}\,\tilde{X}_c \leq \min(M-1,N)$ & $\leq 149$ \\
$K_c$ & retained modes, channel $c$ & 20 \\
$n_{\mathrm{GP}}$ & scalar GPs, $\sum_{c \in \mathcal{C}} K_c$ & 20 \\
\bottomrule
\end{tabular}
\caption{Dimensions of the full-order model, the grid interpolation, the reduced basis and
the regressor.}
\label{tab:dims}
\end{table}

\section{Measured results}
\label{sec:measured}

\subsection{Spectra}

For $M = 150$ realizations, a basis built from $M_{\mathrm{tr}} = 120$ and evaluated on the
$|\mathcal{T}| = 30$ held out, \eqref{eq:error-proj} gives for the mass channel $c=1$ at
terminal time $T = 2$\,s:

\begin{center}
\begin{tabular}{lrrrrr}
\toprule
$K_1$ & 0 (mean field only) & 5 & 10 & 20 & 40 \\
\midrule
$\varepsilon^{\text{proj}}_1$ & 0.871 & 0.516 & 0.381 & 0.243 & 0.192 \\
\bottomrule
\end{tabular}
\end{center}

with $\eta_1 = 0.90$ at $K_1 = 11$ and $\eta_1 = 0.99$ at $K_1 = 51$. A linear basis is
adequate for this channel: $K_1 = 20$ reduces a $16\,384$-dimensional field to a relative
error of $0.24$, against $0.87$ for the mean field alone.

The momentum channels admit no useful representation at $T$, with
$\varepsilon^{\text{proj}}_c(20) \approx 0.98$ for $c = 2,3,4$. This reflects the state and
not the method. The ratio of root-mean-square (RMS) momentum density to RMS mass density
falls from $2.29$ at $t = 0.2$\,s to $0.018$ at $t = 2$\,s: the assembly is at rest and
$\rho \bs{u}$ is residual creep. Evaluated at $t = 0.2$\,s the same construction gives
$\varepsilon^{\text{proj}}_c(20) = 0.35$--$0.40$ for $c = 2,3,4$ against $0.34$ for $c = 1$.
The momentum channels are thus informative during transport and vacuous at rest, consistent
with their role in Section~\ref{sec:encoder}. Since the initial velocity enters
$\bs{\theta}$ directly, no information is lost by taking $\mathcal{C} = \{1\}$ for the
terminal-state surrogate.

\subsection{Composed surrogate}

With $\mathcal{C} = \{1\}$, $K_1 = 20$, $n_{\mathrm{GP}} = 20$, $M_{\mathrm{tr}} = 120$ and
$|\mathcal{T}| = 30$:

\begin{center}
\begin{tabular}{lr}
\toprule
Quantity & Relative error \\
\midrule
Mean field only, $K_1 = 0$ & 0.871 \\
$\varepsilon^{\text{proj}}_1(20)$, truncation only, Eq.~\eqref{eq:error-proj} & 0.243 \\
$\varepsilon^{\text{sur}}_1(20)$, composed surrogate, Eq.~\eqref{eq:error-sur} & 0.381 \\
Regression alone, in coefficient space & 0.350 \\
\bottomrule
\end{tabular}
\end{center}

and $\eta_1(20) = 0.968$. The truncation and regression contributions are comparable, so
neither refining the basis nor improving the regressor alone would reduce
$\varepsilon^{\text{sur}}_1$ substantially. The per-mode coefficient of determination on
$\mathcal{T}$ is $R^2 \geq 0.93$ for $k \leq 5$, decaying to $R^2 \approx 0.73$--$0.86$ by
$k = 9$: the leading modes carry smooth, parameter-driven structure and the tail carries
realization-specific detail.

Pooling all modes, the standardized residuals
$z_{c,k} = (m_{c,k} - a_{c,k}) / v_{c,k}^{1/2}$ satisfy
$\mathrm{Pr}(|z| \leq 1) = 0.743$ and $\mathrm{Pr}(|z| \leq 2) = 0.918$ against nominal
$0.683$ and $0.954$, with $\mathrm{RMS}(z) = 1.16$. The predictive variance
\eqref{eq:gp-var} is close to calibrated and mildly overconfident.

The leading modes admit direct physical interpretation. Over the training set the
correlation between $(\bs{a}_1)_k$ and the components of $\bs{\theta}$ is maximal at $0.673$
for $k=1$ against the first velocity component, $0.776$ for $k=2$ against the second, and
$0.898$ for $k=6$ against the particle mass; the singular values $\sigma_{1,1} = 107.8$ and
$\sigma_{1,2} = 101.7$ are nearly equal, as required by the symmetry of $\Omega$ and
$\mathcal{P}$ under exchange of the first two coordinates.

Ranking inputs by ARD relevance, the $\sigma^2$-weighted mean of $1/\ell_i$ over retained
modes,
\begin{equation}
  v_{x}\,(1.94) \; \approx \; v_{y}\,(1.66) \; > \; \mu\,(1.25) \; > \;
  v_{z}\,(0.90) \; > \; m\,(0.78) \; \approx \; e\,(0.74) \; \gg \; z_0\,(0.23).
\end{equation}
Aggregating $\ell_i$ rather than $1/\ell_i$ is incorrect: the basis assigns the two
horizontal velocity components to separate modes of nearly equal energy, so a weighted mean
of length scales ranks $v_x$ above $v_y$ by the margin
$\sigma_{1,1}^2 - \sigma_{1,2}^2$ alone, an artifact of the aggregation rather than a
property of the system.

\subsection{Irreducible scatter}
\label{sec:floor}

The errors above are uninterpretable without the scatter of the FOM itself at fixed
$\bs{\theta}$. We generate an ensemble of $E$ realizations sharing one $\bs{\theta}$ and
differing only by an independent random displacement of magnitude $\epsilon$ applied to
every initial particle position, and measure separation in two senses: the root-mean-square
(RMS) separation of corresponding particles,
\begin{equation}
  D_{\text{grain}}(t) = \left( \frac{1}{n_p} \sum_{i=1}^{n_p}
  \norm{ \bs{x}_i^{(e)}(t) - \bs{x}_i^{(0)}(t) }_2^2 \right)^{1/2},
  \label{eq:grain-divergence}
\end{equation}
and the relative difference of the mass field,
$D_{\text{field}}(t) = \norm{\bs{s}_1^{(e)}(t) - \bs{s}_1^{(0)}(t)}_2 /
\norm{\bs{s}_1^{(0)}(t)}_2$.

For $\epsilon = 10^{-3} r$, with $r$ the particle radius, $D_{\text{grain}}$ amplifies by a
factor $3.8 \times 10^{3}$, from $10^{-5}$\,m to $3.81 \times 10^{-2}$\,m, the latter being
$1.9$ particle diameters. The growth is not a single exponential: the local rate
$\mathrm{d}(\ln D_{\text{grain}})/\mathrm{d}t$ is $23.3\,\mathrm{s}^{-1}$ on
$t \in [0.02, 0.06]$\,s, falls to a sustained $5.6$--$6.4\,\mathrm{s}^{-1}$ on
$[0.14, 0.60]$\,s, corresponding to an $e$-folding time
$\lambda^{-1} = 0.16$--$0.18$\,s, and collapses to $0.17\,\mathrm{s}^{-1}$ thereafter. The
plateau is reached, to within ten percent, at $t = 0.62$\,s, which we take as the horizon
over which individual positions remain predictable.

The plateau is set by jamming rather than by the assembly scale: growth ceases when the pile
comes to rest and frictional contacts lock the packing, which is why $D_{\text{grain}}$
saturates at $1.9$ diameters and not at the extent of the domain. The estimate
$\lambda^{-1} \ln (D_{\text{sat}}/\epsilon) = 1.38$\,s consequently overstates the measured
horizon, the early growth being faster than the sustained rate.

Over the same interval $D_{\text{field}}$ reaches only $0.086$ and the displacement of the
field center of mass $0.009$ grid nodes. Chaos destroys the trajectory and leaves the
coarse-grained functional intact.

For the surrogate the relevant quantity is the scatter of realizations about their
conditional mean, not the separation of two realizations. If
$\bs{s}_1 = \bs{\mu} + \bs{n}$ with $\E[\bs{n}] = \bs{0}$ and $\bs{n}$ independent between
realizations, then
\begin{equation}
  \E \norm{\bs{s}_1^{(a)} - \bs{s}_1^{(b)}}_2^2 = 2\,\E\norm{\bs{n}}_2^2 ,
  \label{eq:sqrt2}
\end{equation}
so a model predicting $\bs{\mu}$ incurs $\norm{\bs{n}}$ while the pairwise difference is
larger by $\sqrt{2}$. With $\epsilon = 0.25 r$, comparable to the lattice disorder, and
$E = 16$, the measured pairwise difference is $0.245$ and the measured RMS deviation about
the ensemble mean is
\begin{equation}
  \varepsilon^{\text{floor}}_1 = \left(
  \frac{\sum_{e} \norm{ \bs{s}_1^{(e)} - \bar{\bs{s}}_1^{\,\text{ens}} }_2^2}{
  \sum_{e} \norm{ \bs{s}_1^{(e)} }_2^2} \right)^{1/2} = 0.168 ,
  \label{eq:floor}
\end{equation}
their ratio being $1.46$ against the $\sqrt{2} = 1.414$ predicted by \eqref{eq:sqrt2}.

Relative to \eqref{eq:floor}, $\varepsilon^{\text{proj}}_1(20) = 0.243$ is $1.45$ times the
floor and $\varepsilon^{\text{sur}}_1(20) = 0.381$ is $2.27$ times it. Both exceed the floor,
so both admit improvement; substituting the pairwise $0.245$ for \eqref{eq:floor} would have
indicated, incorrectly, that the surrogate was near-optimal. Since
$\varepsilon^{\text{proj}}_1(40) = 0.192$ remains above $0.168$, the truncation rank is not
yet sufficient at $K_1 = 20$, and the gap
$\varepsilon^{\text{sur}}_1 - \varepsilon^{\text{proj}}_1$ indicates the regression as the
larger of the two deficits at $M_{\mathrm{tr}} = 120$.

\section{Latent dynamics}
\label{sec:dynamics}

\subsection{Formulation}

Section~\ref{sec:gp} regresses the terminal state directly from $\bs{\theta}$. An alternative is to
learn a one-step map in the reduced coordinates and apply it repeatedly. Partition $\bs{\theta}$ as
$\bs{\theta} = (\bs{\theta}_a, \bs{\theta}_m)$, with $\bs{\theta}_a \in \R^{4}$ the action
components (the three initial velocity components and the release height) and
$\bs{\theta}_m \in \R^{3}$ the material components (particle mass, restitution, friction). Let
\begin{equation}
  \bs{z}(t) = \big[\, \bs{a}_1(t)^{\top}, \ldots, \bs{a}_{n_c}(t)^{\top} \,\big]^{\top}
  \in \R^{n_z}, \qquad n_z = \sum_{c=1}^{n_c} K_c ,
  \label{eq:latent-state}
\end{equation}
be the concatenated modal coordinates of all $n_c$ channels, each obtained by
\eqref{eq:projection}. We regress
\begin{equation}
  \mathcal{G} : \R^{n_z + 3} \to \R^{n_z},
  \qquad
  \bs{z}(t + \Delta) = \mathcal{G}\big( \bs{z}(t), \bs{\theta}_m \big),
  \label{eq:step-map}
\end{equation}
component-wise by $n_z$ independent Gaussian processes as in Section~\ref{sec:gp}, and define the
$j$-step prediction by composition,
\begin{equation}
  \hat{\bs{z}}(t_0 + j\Delta) = \mathcal{G}^{\,j}\big( \bs{z}(t_0), \bs{\theta}_m \big),
  \label{eq:rollout}
\end{equation}
with $\bs{z}(t_0)$ the projection of an observed initial field. Here $\Delta = 0.1$\,s and the solver
timestep is $10^{-4}$\,s, so each application of $\mathcal{G}$ advances $10^{3}$ solver steps.

Two choices are consequential. First, all $n_c = 4$ channels are carried in \eqref{eq:latent-state},
whereas Section~\ref{sec:measured} predicts only $c=1$: the mass field alone is not a Markov state,
since two assemblies with equal mass fields and unequal velocity fields evolve differently. Second,
$\bs{\theta}_a$ is deliberately \emph{withheld} from \eqref{eq:step-map}. Supplying it would permit
$\mathcal{G}$ to ignore $\bs{z}$ and internally represent $\bs{\theta} \mapsto \bs{z}(t+\Delta)$,
collapsing the construction onto the direct regression of Section~\ref{sec:gp} and rendering the
sufficiency of $\bs{z}$ untestable.

\subsection{Error decomposition}

The error of \eqref{eq:rollout} separates into three terms, each independently measurable, and the
identity of the dominant term determines which remedy is appropriate. Writing
$\varepsilon^{\text{proj}}$ for \eqref{eq:error-proj} evaluated at time $t$,
$\varepsilon^{\text{step}}$ for the error of a single application of $\mathcal{G}$ to the
\emph{projected} state (teacher forcing), and $\varepsilon^{\text{roll}}$ for that of
\eqref{eq:rollout},
\begin{equation}
  \underbrace{\varepsilon^{\text{proj}}}_{\text{truncation}}
  \; \leq \;
  \underbrace{\varepsilon^{\text{step}}}_{+\;\text{regression}}
  \; \leq \;
  \underbrace{\varepsilon^{\text{roll}}}_{+\;\text{accumulation}} .
  \label{eq:three-terms}
\end{equation}
At $K_1 = 16$ we measured $\varepsilon^{\text{proj}} = 0.4012$ and
$\varepsilon^{\text{step}} = 0.4013$ at $t = 1$\,s: the regression term is $10^{-4}$, so the
truncation term accounts for $86\%$ of $\varepsilon^{\text{roll}} = 0.464$ and no quantity of
additional training data could have improved the result. Raising the rank to $K_1 = 96$ reduced
$\varepsilon^{\text{proj}}$ to $0.173$ and left accumulation dominant, at which point reducing the
number of compositions in \eqref{eq:rollout} and increasing the density of the training design became
the appropriate remedies.

\subsection{Results}

With $K_1 = 96$, $K_{2,3,4} = 24$ so that $n_z = 168$, $\Delta = 0.1$\,s, $M = 1536$ realizations of
$N_t = 10$ instants, and $M_{\mathrm{tr}} = 900$ transitions drawn from the $12\,744$ available:

\begin{center}
\begin{tabular}{lrr}
\toprule
Quantity at $t = 1$\,s & $\Delta = 0.05$\,s, $M = 256$ & $\Delta = 0.10$\,s, $M = 1536$ \\
\midrule
$\varepsilon^{\text{proj}}$ (truncation) & 0.173 & 0.108 \\
$\varepsilon^{\text{step}}$ (teacher forced) & 0.177 & 0.109 \\
$\varepsilon^{\text{roll}}$ (rollout) & 0.445 & \textbf{0.174} \\
persistence, $\bs{z}(t+\Delta) = \bs{z}(t)$ & 1.762 & 1.531 \\
accumulation, $\varepsilon^{\text{roll}} - \varepsilon^{\text{step}}$ & 0.268 & 0.065 \\
$\varepsilon^{\text{floor}}$ for that configuration & 0.168 & 0.154 \\
$\varepsilon^{\text{roll}} / \varepsilon^{\text{floor}}$ & 2.65 & \textbf{1.13} \\
\bottomrule
\end{tabular}
\end{center}

$\varepsilon^{\text{roll}}$ is non-monotonic in $t$, attaining $0.33$ at $t \approx 0.35$\,s before
decreasing: the impact and spreading phase is the least predictable part of the trajectory, after
which the assembly jams and the state becomes easier to predict.

The improvement was obtained by halving the number of compositions in \eqref{eq:rollout} and by
contracting $\mathcal{P}$ by approximately $40\%$ along each axis. The latter reduces
$\varepsilon^{\text{step}}$ at fixed $M_{\mathrm{tr}}$, since a fixed training set covers a smaller
domain more densely, and simultaneously reduces $\varepsilon^{\text{proj}}$, since a less diverse
ensemble is spanned by fewer modes ($\eta_1(96)$ rising from $0.9943$ to $0.9974$). It is
purchased at the cost of generality: the contracted $\mathcal{P}$ has approximately $7\%$ of the
original volume in $p = 7$ dimensions, and $\varepsilon^{\text{roll}} = 0.174$ is meaningful only on
the contracted domain.

We note that exact Gaussian process inference is $O(M_{\mathrm{tr}}^3)$ per evaluation of
\eqref{eq:mll}, so $M_{\mathrm{tr}}$ cannot be raised in proportion to $M$. A larger $M$ improves
the basis and permits a better-covering subset, but reducing $\varepsilon^{\text{step}}$ at fixed
$M_{\mathrm{tr}}$ requires contracting $\mathcal{P}$ rather than enlarging the design.

\section{The irreducible floor}
\label{sec:floor-theory}

Every error reported above is quoted against $\varepsilon^{\text{floor}}$, and the definition of that
quantity, its justification, and the manner of its measurement are given here in full, since a
surrogate error is not interpretable without it.

\subsection{Definition}

Let $\bs{\theta}$ be fixed and let $\omega$ denote the microstate, that is, the particular initial
packing of the assembly. The vector $\bs{\theta}$ specifies the \emph{macrostate}; it does not
determine $\omega$, and many $\omega$ are consistent with one $\bs{\theta}$. The terminal field is
therefore a random variable $\bs{s}_1(\bs{\theta}, \omega)$, and any surrogate
$\hat{\bs{s}}_1(\bs{\theta})$ is a function of $\bs{\theta}$ alone. Write
\begin{equation}
  \bs{\mu}(\bs{\theta}) = \E_\omega \big[ \bs{s}_1(\bs{\theta}, \omega) \big],
  \qquad
  \bs{n}(\bs{\theta}, \omega) = \bs{s}_1(\bs{\theta}, \omega) - \bs{\mu}(\bs{\theta}),
  \qquad \E_\omega[\bs{n}] = \bs{0}.
  \label{eq:decomposition}
\end{equation}

\subsection{Justification}

Expanding the expected squared error of an arbitrary $\hat{\bs{s}}_1$ and using
$\E_\omega[\bs{n}] = \bs{0}$ to annihilate the cross term,
\begin{equation}
  \E_\omega \norm{ \hat{\bs{s}}_1(\bs{\theta}) - \bs{s}_1(\bs{\theta}, \omega) }_2^2
  = \underbrace{\norm{ \hat{\bs{s}}_1(\bs{\theta}) - \bs{\mu}(\bs{\theta}) }_2^2}_{\text{model error}}
  + \underbrace{\E_\omega \norm{ \bs{n}(\bs{\theta}, \omega) }_2^2}_{\text{irreducible}} .
  \label{eq:bias-variance}
\end{equation}
The second term does not contain $\hat{\bs{s}}_1$. It is therefore unaffected by any change to the
surrogate, and the first term is minimised, at zero, by $\hat{\bs{s}}_1 = \bs{\mu}$. The best
attainable error is thus $\big( \E_\omega \norm{\bs{n}}_2^2 \big)^{1/2}$, and it is unattainable to
improve upon because the information distinguishing realizations is absent from the input, not merely
unexploited. In the terminology of uncertainty quantification this is the \emph{aleatoric} component,
as against the \emph{epistemic} component represented by the first term.

\subsection{Measurement}

We generate $E$ realizations sharing $\bs{\theta}$ and differing only by an independent random
displacement of magnitude $\epsilon$ applied to each initial particle position, take $\epsilon$
comparable to the disorder of the packing itself so that the realizations are genuinely distinct
members of one macrostate, and evaluate the normalized scatter about their sample mean,
\begin{equation}
  \varepsilon^{\text{floor}}
  = \left( \frac{\sum_{e=1}^{E} \norm{ \bs{s}_1^{(e)} - \bar{\bs{s}}_1^{\,\text{ens}} }_2^2}
  {\sum_{e=1}^{E} \norm{ \bs{s}_1^{(e)} }_2^2} \right)^{1/2},
  \qquad
  \bar{\bs{s}}_1^{\,\text{ens}} = \frac{1}{E} \sum_{e=1}^{E} \bs{s}_1^{(e)} .
  \label{eq:floor-measured}
\end{equation}

Three properties of \eqref{eq:floor-measured} require comment.

\emph{It is not the pairwise separation.} Substituting the pairwise difference between two
realizations overstates the floor by $\sqrt{2}$: from \eqref{eq:decomposition} with $\bs{n}$
independent between realizations,
\begin{equation}
  \E \norm{ \bs{s}_1^{(a)} - \bs{s}_1^{(b)} }_2^2
  = \E \norm{ \bs{n}^{(a)} - \bs{n}^{(b)} }_2^2
  = 2\, \E \norm{ \bs{n} }_2^2 ,
  \label{eq:sqrt2-again}
\end{equation}
whereas a model predicting $\bs{\mu}$ incurs $\E\norm{\bs{n}}_2^2$. Measured on the contracted
configuration at $t = 1$\,s with $E = 16$: \eqref{eq:floor-measured} gives $0.154$ while the pairwise
value is $0.225$, a ratio of $1.463$ against the $\sqrt{2} = 1.414$ predicted by
\eqref{eq:sqrt2-again}. Quoting the pairwise figure would understate
$\varepsilon^{\text{roll}}/\varepsilon^{\text{floor}}$ from $1.13$ to $0.77$, that is, would report a
surrogate as having surpassed a bound it had not reached.

\emph{It is configuration-specific.} $\varepsilon^{\text{floor}}$ depends on $\bs{\theta}$, on the
horizon, and on the grid; it is $0.168$ for the wide-$\mathcal{P}$ configuration at $t = 2$\,s and
$0.154$ for the contracted configuration at $t = 1$\,s. It must be re-measured whenever any of these
changes.

\emph{The sample mean biases it low.} Using $\bar{\bs{s}}_1^{\,\text{ens}}$ in place of the unknown
$\bs{\mu}$ underestimates the scatter by the factor $\sqrt{(E-1)/E}$, which is $0.968$ at $E = 16$.
The reported figures are not corrected for this, and it is small relative to the variation of
$\varepsilon^{\text{floor}}$ with $\bs{\theta}$, which has not been characterized: all values here
are measured at a single, central $\bs{\theta}$.

\emph{Finally, the floor is not a reconstruction bound.} $\varepsilon^{\text{proj}} = 0.108$ lies
\emph{below} $\varepsilon^{\text{floor}} = 0.154$, and this is not contradictory: the former is the
error of compressing a field that has already been realized, the latter the error of predicting which
field will be realized. Only quantities of the second kind are bounded by
\eqref{eq:floor-measured}.

\section{Decoding to particles}
\label{sec:decode}

\subsection{The decode is a sampler}

The interpolation $\mathcal{E}$ of Section~\ref{sec:encoder} is many-to-one: distinct particle
configurations produce identical fields, since \eqref{eq:deposit} retains only $n_c N$ numbers of the
$n_q$ that describe $\bs{q}$, and here $n_c N / n_q = 3.2$ but the map is not injective because
positions within a cell are not resolved. The preimage $\mathcal{E}^{-1}(\bs{s})$ is therefore a set,
and no inverse exists. We construct instead a sampler
\begin{equation}
  \mathcal{D} : \R^{n_c N} \to \R^{n_q},
  \qquad
  \mathcal{D}(\bs{s}) \sim \pi\big( \cdot \mid \bs{s} \big),
  \label{eq:decoder}
\end{equation}
drawing one element of that set. Consequently $\mathcal{D}(\mathcal{E}(\bs{q})) \neq \bs{q}$: the
construction reproduces the field, not the configuration.

\subsection{Construction}

Given a predicted mass field $\bs{\rho}$ and momentum fields $(\bs{\rho u})_j$, and given that the
per-particle mass $m$ and hence the total mass $M = n_p m$ are known from $\bs{\theta}$:

\emph{(i) Non-negativity.} A truncated expansion $\bs{\rho} = \bar{\bs{s}}_1 + \Phi_1 \bs{a}_1$ is not
constrained to be non-negative and is not, so set $\tilde{\rho}_\alpha = \max(\rho_\alpha, 0)$. The
discarded fraction $\sum_\alpha \max(-\rho_\alpha,0) / \sum_\alpha \max(\rho_\alpha,0)$ is $0.04$ to
$0.10$ in the cases reported below.

\emph{(ii) Thresholding.} Truncation also produces small positive values far from the material. Set
$\tilde{\rho}_\alpha \leftarrow 0$ wherever $\tilde{\rho}_\alpha < \tau \max_\beta
\tilde{\rho}_\beta$, with $\tau = 0.02$. This step is not cosmetic. Individually negligible, those
values are spread over all $N = 16\,384$ nodes and collectively attract a substantial number of
particles into a halo absent from the truth; and because they sit where $\tilde{\rho} \approx 0$, the
quotient in \eqref{eq:decode-velocity} assigns those particles very large velocities. Measured at
$t = 0.7$\,s, thresholding reduces the error in the sampled spatial standard deviation from
$0.092$\,m to $0.006$\,m and the sampled kinetic energy from $141$ to $7.6$ against a true value of
$3.6$.

\emph{(iii) Mass normalisation.} Set $\hat{\rho}_\alpha = \tilde{\rho}_\alpha M / \sum_\beta
\tilde{\rho}_\beta$, so that the decode conserves mass exactly even though steps (i) and (ii) do not.

\emph{(iv) Positions.} Two samplers are used. The first draws a node index
$\alpha \sim \hat{\rho}_\alpha / M$ and then a position uniformly in that node's cell; it reproduces
the nodal mass in expectation and ignores the packing constraint. The second is described in
Section~\ref{sec:poisson}.

\emph{(v) Velocities.} With $w_\alpha(\bs{x})$ the same weights \eqref{eq:weights} used to deposit,
\begin{equation}
  \bs{v}(\bs{x}) = \frac{\sum_\alpha w_\alpha(\bs{x}) \, (\bs{\rho u})_\alpha}
  {\sum_\alpha w_\alpha(\bs{x}) \, \hat{\rho}_\alpha},
  \label{eq:decode-velocity}
\end{equation}
evaluated at the sampled position and set to zero where the denominator vanishes. Using the deposit's
own weights rather than the nodal value keeps $\bs{v}$ continuous across cell boundaries.

\subsection{Minimum-separation sampling}
\label{sec:poisson}

The uniform sampler of (iv) produces a configuration in which particles overlap: node spacing is
$h_1 = h_2 = 0.065$\,m against a diameter $2r = 0.02$\,m, so roughly $30$ particles share a node and
uniform placement within it is Poisson-like rather than packed. Such a configuration is unusable as an
initial condition, for the reason quantified in Section~\ref{sec:restart}.

We therefore impose
\begin{equation}
  \norm{ \bs{x}_i - \bs{x}_j }_2 \; \geq \; 2r \quad \forall\, i \neq j,
  \qquad
  \Omega_j^{\min} + r \leq (\bs{x}_i)_j \leq \Omega_j^{\max} - r ,
  \label{eq:poisson-constraint}
\end{equation}
the second condition ensuring a particle centre is a full radius clear of each wall. Candidates are
proposed as in (iv) and accepted only if \eqref{eq:poisson-constraint} holds, with a spatial hash of
cell size $2r$ so that a candidate need be tested against the $3^d = 27$ neighbouring cells only.
Bridson's algorithm is the usual construction for a Poisson-disk set but generates a spatially
\emph{uniform} intensity; here the intensity is prescribed by $\hat{\rho}$ and the disk radius is
fixed by the physics, so the density is carried by the proposal distribution instead.

Candidates violating the second condition of \eqref{eq:poisson-constraint} must be \emph{rejected} and
not projected onto the boundary. Projection collapses all such candidates onto a face --- and since the
assembly rests on $\Omega_3^{\min}$, a large fraction of proposals are of this kind --- which destroys
the separation just enforced; we measured $36\%$ of pairs in violation of the first condition despite a
correct acceptance test.

The construction cannot produce a density exceeding the maximum packing fraction. Where $\hat{\rho}$
demands one, the requested particles cannot be placed; this indicates that the predicted field is
locally unphysical, and the deficit is reported rather than concealed. In the cases below no deficit
occurred: all $n_p = 3375$ particles were placed.

Measured on the resulting configuration: $\min_{i \neq j} \norm{\bs{x}_i - \bs{x}_j}_2 / 2r = 1.0002$,
median nearest-neighbour distance $1.157 \times 2r$, and $0\%$ of pairs in contact, against $73\%$ for
the uniform sampler at a median of $0.699 \times 2r$.

\subsection{Restart validation}
\label{sec:restart}

Whether $\mathcal{D}(\bs{s})$ is admissible as an initial condition is decided by integrating from it.
Both configurations were written in the solver's state layout and advanced; the total energy of
Section~\ref{sec:pod} (kinetic, gravitational and contact-elastic) at $t=0$ and thereafter is

\begin{center}
\begin{tabular}{lrrl}
\toprule
Sampler & pairs in contact & $E(0)$ & behaviour \\
\midrule
uniform & $73\%$ & $14\,893.5$ & $90\%$ of $E$ released within $0.01$\,s \\
minimum-separation & $0\%$ & $1\,270.8$ & monotone decay \\
\bottomrule
\end{tabular}
\end{center}

The factor of $11.7$ is elastic energy stored in overlapping contacts by \eqref{eq:deposit}'s inverse
being taken carelessly, and it is converted immediately into kinetic energy. From the
minimum-separation configuration the energy decreases monotonically over the full second integrated,
$1270.8 \to 78.3$, with no contact-list overflow: the state is admissible.

\subsection{Limitations}

Two are intrinsic to the channels retained.

Only the first velocity moment is stored. The within-node velocity distribution --- the granular
temperature --- is discarded by \eqref{eq:deposit}, so \eqref{eq:decode-velocity} assigns every
particle in a neighbourhood the same velocity and the sampled kinetic energy is biased. Retaining a
second-moment channel $\rho u_i u_j$ would recover it.

Enforcing \eqref{eq:poisson-constraint} costs fidelity in the low-order moments, because the sampler
cannot place particles at a density above the packing limit and displaces them outward instead. At
$t = 1$\,s the error in the centroid rises from $0.025$\,m to $0.057$\,m and in the spatial standard
deviation from $0.011$\,m to $0.045$\,m, against a particle diameter of $0.02$\,m. The uniform sampler
is therefore preferable for visualisation and for bulk statistics, and the minimum-separation sampler
whenever the configuration is to be integrated.

\section*{Nomenclature}

\begin{tabbing}
\hspace{2.6cm} \= \kill
ARD \> automatic relevance determination (one kernel length scale per input) \\
FOM \> full-order model (the particle dynamics simulation) \\
GP \> Gaussian process \\
POD \> proper orthogonal decomposition \\
QMC \> quasi-Monte Carlo \\
RMS \> root mean square \\
SVD \> singular value decomposition \\
\end{tabbing}

\end{document}
